\section*{Capítulo \ref{ch:b1:heat-budget} --- O balanço térmico da Terra}

\begin{solution}{exo:b1:heat-budget:1}
$q = k\,\mathrm dT/\mathrm dz = 2.5\times\qty{0.022}{K.m^{-1}} =
\qty{0.055}{W.m^{-2}} = \qty{55}{mW.m^{-2}}$. Através de \qty{1}{km^2} $=
\qty{e6}{m^2}$: \qty{5.5e4}{W}, isto é, \qty{55}{kW}.
\end{solution}

\begin{solution}{exo:b1:heat-budget:2}
Continentes $\num{2.073e14}\times0.0709 = \qty{1.470e13}{W} = \qty{14.70}{TW}$;
oceanos $\num{3.028e14}\times0.1054 = \qty{31.92}{TW}$; total \qty{46.61}{TW}, e
o fluxo médio de calor $\num{4.661e13}/\num{5.101e14} = \qty{91.4}{mW.m^{-2}}$ (o
valor publicado, 46.7 e 91.6, difere apenas pelo arredondamento das duas médias). A
fração da entrada solar é $46.6/173000 = \num{2.7e-4}$, ou
\qty{0.027}{\percent}.
\end{solution}

\begin{solution}{exo:b1:heat-budget:3}
Urânio: $0.99274\times\num{9.513e-5} + 0.00720\times\num{5.685e-4} =
\num{9.444e-5} + \num{4.09e-6} = \qty{9.85e-5}{W.kg^{-1}}$; o raro
\ce{^{235}U} fornece \qty{4}{\percent} dele. Potássio: a fração mássica de
\ce{^{40}K} é $\num{1.17e-4}\times39.96/39.10 = \num{1.196e-4}$, logo
$\num{1.196e-4}\times\num{2.847e-5} = \qty{3.40e-9}{W.kg^{-1}}$.
\end{solution}

\begin{solution}{exo:b1:heat-budget:4}
A conversão é $\qty{1}{\micro cal.cm^{-2}.s^{-1}} = \num{4.184e-6}\,\unit{J.s^{-1}}/\qty{e-4}{m^2}
= \qty{41.84}{mW.m^{-2}}$: leste dos Estados Unidos \qty{33.5}{mW.m^{-2}}, Basin and
Range \qty{58.6}{mW.m^{-2}}, Sierra Nevada \qty{16.7}{mW.m^{-2}}. Com $A_0D$ em
\unit{mW.m^{-2}} quando $A_0$ está em \unit{\micro W.m^{-3}} e $D$ em \unit{km}:
$33.5 + 3.0\times7.5 = \qty{56.0}{mW.m^{-2}}$, $58.6 + 3.0\times9.4 =
\qty{86.8}{mW.m^{-2}}$, $16.7 + 3.0\times10 = \qty{46.7}{mW.m^{-2}}$. O mesmo
granito dá \omterm{def:b1:heat-budget:heat-flow}{fluxos de calor} quase duas vezes maiores na Basin and Range do que na
Sierra Nevada: a diferença está no \omterm{def:b1:layered-earth:layers}{manto} abaixo.
\end{solution}

\begin{solution}{exo:b1:heat-budget:5}
$T(10\,\text{km}) = 10 + \frac{0.060\times10^4}{2.5} - \frac{\num{2e-6}\times10^8}{2\times2.5}
= 10 + 240 - 40 = \qty{210}{\degreeCelsius}$. Abaixo, $q_m = 60 - 2.0\times10 =
\qty{40}{mW.m^{-2}}$, um gradiente de \qty{16}{K.km^{-1}}, e
$T(35\,\text{km}) = 210 + 16\times25 = \qty{610}{\degreeCelsius}$.
\end{solution}

\begin{solution}{exo:b1:heat-budget:6}
Folhelho: gradiente $12.0/0.300 = \qty{40}{K.km^{-1}}$, $q = 1.5\times0.040 =
\qty{60}{mW.m^{-2}}$. Granito: $10.0/0.500 = \qty{20}{K.km^{-1}}$,
$q = 3.0\times0.020 = \qty{60}{mW.m^{-2}}$. $R(300) = 300/1.5 =
\qty{200}{m^2.K.W^{-1}}$, $R(800) = 200 + 500/3.0 = \qty{366.7}{m^2.K.W^{-1}}$.
Os três pontos $(0, 12.0)$, $(200, 24.0)$, $(366.7, 34.0)$ estão alinhados, com
inclinação $22.0/366.7 = \qty{0.0600}{W.m^{-2}}$: o \omterm{def:b1:heat-budget:heat-flow}{fluxo de calor} é o mesmo,
\qty{60}{mW.m^{-2}}, nas duas rochas; só o gradiente muda, em proporção inversa
à condutividade.
\end{solution}

\begin{solution}{exo:b1:heat-budget:7}
$\tau = L^2/\kappa$: \qty{1}{m}, \qty{e6}{s} ou \num{12}~dias; \qty{10}{m},
\qty{e8}{s} ou \num{3.2}~anos; \qty{100}{km}, \qty{e16}{s} ou
\qty{317}{Myr}; \qty{2900}{km}, \qty{8.4e18}{s} ou \qty{267}{Gyr}. A
\qty{10}{m} o ciclo diário não é sentido de forma alguma, e o ciclo anual, mais curto
que os três anos necessários para atravessar \qty{10}{m}, só muito fracamente: passam
apenas as mudanças que duram anos ou mais. O último caso dá $267/4.54 = 59$ vezes a
idade da Terra.
\end{solution}

\begin{solution}{exo:b1:heat-budget:8}
(a) $1646 - 0.54\times50 = \qty{1619}{K}$, isto é \qty{1346}{\degreeCelsius},
o valor publicado de \qty{1350}{\degreeCelsius} arredondado. (b) Gradiente médio
$(0.54 + 0.36)/2 = \qty{0.45}{K.km^{-1}}$, logo $1646 + 0.45\times360 =
\qty{1808}{K}$, \qty{9}{K} (\qty{0.5}{\percent}) acima de \qty{1799}{K}: o
gradiente não diminui linearmente, mas a estimativa está próxima.
\end{solution}

\begin{solution}{exo:b1:heat-budget:9}
(a) $G^2 = T_0^2/(\pi\kappa t)$ dá $t = T_0^2/(\pi\kappa G^2)$.
(b) $G = 0.02\times\tfrac59\,\unit{K}/\qty{0.3048}{m} = \qty{0.03645}{K.m^{-1}} =
\qty{36.45}{K.km^{-1}}$; $t = 3900^2/(\pi\times\num{1.2e-6}\times0.03645^2) =
\num{1.521e7}/\num{5.009e-9} = \qty{3.04e15}{s}$, isto é \qty{96}{Myr},
a figura de Kelvin. (c) $t = 1300^2/(\pi\times\num{e-6}\times0.036^2) =
\qty{4.15e14}{s} = \qty{13}{Myr}$: uma temperatura inicial mais realista torna a
Terra condutiva ainda mais jovem.
\end{solution}

\begin{solution}{exo:b1:heat-budget:10}
(a) Depois da erosão, a nova superfície é a antiga profundidade $e$, onde a
produção de calor era $A_0' = A_0e^{-e/D}$ e o \omterm{def:b1:heat-budget:heat-flow}{fluxo de calor} $q(e) = q_m +
A_0De^{-e/D} = q_m + A_0'D$; medida a partir da nova superfície, a produção ainda
decresce como $A_0'e^{-z'/D}$. O plúton desce ao longo da mesma reta.
(b) A camada agora tem espessura $h - e$ com o mesmo $A$, logo $q_0' = q_m + A(h - e)$:
sobre uma reta de inclinação $h - e$ passando por $q_m$. (c) Só o modelo exponencial: no
modelo de camada uniforme, plútons erodidos a profundidades diferentes estariam sobre retas de
inclinações diferentes, e os pontos de uma província se espalhariam. Que uma única
reta seja observada favorece uma produção que decresce aproximadamente exponencialmente com a
profundidade.
\end{solution}

\begin{solution}{exo:b1:heat-budget:11}
(a) $\tau_P = RL/(3\kappa)$: para \qty{50}{km}, $\num{6.4e6}\times\num{5e4}/\num{3e-6}
= \qty{1.07e17}{s} = \qty{3.4}{Gyr}$; para \qty{100}{km}, \qty{6.8}{Gyr}.
(b) $q = kT_i/L = 3\times1300/\num{5e4} = \qty{78}{mW.m^{-2}}$ no início,
e $78\,e^{-2/3.38} = \qty{43}{mW.m^{-2}}$ após \qty{2}{Gyr}. (c)
$3\times1300/\sqrt{\pi\times\num{e-6}\times\num{6.31e16}} = 3900/\num{4.45e5} =
\qty{8.8}{mW.m^{-2}}$: a Terra bem misturada ainda perde cinco vezes mais calor
após dois bilhões de anos. Todo o seu interior alimenta a superfície, onde numa Terra
condutiva só uma casca externa o faz.
\end{solution}

\begin{solution}{exo:b1:heat-budget:12}
(a) $\sqrt{\pi\kappa\times\qty{1}{Myr}} = \sqrt{\pi\times\num{e-6}\times\num{3.156e13}}
= \qty{9957}{m}$, logo $C = 3\times1336/9957 = \qty{0.403}{W.m^{-2}}$ para
$t$ em Myr, \qty{403}{mW.m^{-2}.Myr^{1/2}}: \qty{82}{\percent} do valor ajustado de 490.
O modelo precisa de um $k\,\Delta T$ maior (cerca de \qty{4900}{W.m^{-1}}): uma
condutividade mais alta do \omterm{def:b1:layered-earth:layers}{manto}, ou um interior mais quente, como ao longo da adiabática.
(b) $490/\sqrt1 = 490$, $490/\sqrt{10} = 155$ e $490/\sqrt{100} =
\qty{49}{mW.m^{-2}}$. (c) $z = 2.326\sqrt{\num{e-6}\times\num{3.156e15}} =
2.326\times\qty{56180}{m} = \qty{131}{km}$.
\end{solution}

\begin{solution}{pb:b1:heat-budget:1}
\textbf{1.} Gradiente $= (T_{500} - T_{100})/\qty{0.4}{km}$: A \qty{15.25}{K.km^{-1}},
B \num{19.50}, C \num{19.50}, D \num{25.50}, E \qty{28.25}{K.km^{-1}}.

\textbf{2.} A perfuração perturba a temperatura da rocha (o fluido de perfuração
esfria ou aquece as paredes), e a perturbação leva semanas a meses para se dissipar.
As dezenas de metros superiores seguem o ciclo sazonal e as mudanças recentes da
temperatura da superfície.

\textbf{3.} $q_0 = k\times$gradiente: A $2.9\times15.25 = \qty{44.225}{mW.m^{-2}}$,
B \qty{50.70}{mW.m^{-2}}, C \qty{60.45}{mW.m^{-2}}, D \qty{68.85}{mW.m^{-2}},
E \qty{84.75}{mW.m^{-2}}.

\textbf{4.} A a D ficam abaixo da média continental, E acima dela em
\qty{13.9}{mW.m^{-2}}, ou \qty{20}{\percent}.

\textbf{5.} O \omterm{def:b1:heat-budget:heat-flow}{fluxo de calor} é condutividade vezes gradiente: C conduz melhor (3.1
contra 2.6), então o mesmo gradiente carrega \qty{19}{\percent} mais calor.

\textbf{6.} As isotermas seguem a topografia só em parte: sob um cume elas sobem menos que a
superfície, logo ficam mais afastadas embaixo dele, e o gradiente medido diminui
(aumenta sob um vale). É preciso uma correção para a topografia antes de comparar os plútons.

\textbf{7.} $\rho C_U H_U$: A $2650\times\num{2.0e-6}\times\num{9.85e-5} =
\qty{0.522}{\micro W.m^{-3}}$, B 0.914, C 1.305, D 1.827, E 2.610.

\textbf{8.} Tório: A 0.558, B 0.976, C 1.255, D 1.742, E 2.439.

\textbf{9.} Potássio: A $2650\times0.030\times\num{3.40e-9} = 0.270$, B 0.315,
C 0.360, D 0.378, E 0.405.

\textbf{10.} $A_0$ = 1.35, 2.20, 2.92, 3.95 e \qty{5.46}{\micro W.m^{-3}}.
O potássio contribui menos, de \qty{20}{\percent} do total em A a
\qty{7}{\percent} em E; urânio e tório dividem o resto quase igualmente.

\textbf{11.} Sim: os cinco pontos sobem regularmente, a cerca de
\qty{1.5}{mW.m^{-2}} de uma reta.

\textbf{12.} $\bar A_0 = 15.88/5 = 3.176$, $\bar q_0 = 308.975/5 = 61.795$;
$\sum(A_0 - \bar A_0)^2 = 10.1681$; $\sum(A_0 - \bar A_0)(q_0 - \bar q_0) =
101.146$.

\textbf{13.} $D = 101.146/10.1681 = 9.95$ (em \unit{mW.m^{-2}} por
\unit{\micro W.m^{-3}}, isto é, em km): $D = \qty{9.95}{km}$, e $q_m = 61.795 -
9.95\times3.176 = \qty{30.20}{mW.m^{-2}}$.

\textbf{14.} A inclinação, cerca de \qty{10}{km}, é a das três províncias; o \omterm{def:b1:heat-budget:heat-flow}{fluxo
de calor} reduzido, \qty{30}{mW.m^{-2}}, está perto dos \qty{33}{mW.m^{-2}} do leste dos
Estados Unidos: uma província calma, com um \omterm{def:b1:layered-earth:layers}{manto} frio embaixo, ao contrário da
Basin and Range.

\textbf{15.} $q_m/k = \qty{12}{K.km^{-1}}$ e $A_0D^2/k = \num{5.46e-6}\times\num{e8}/2.5
= \qty{218.4}{K}$: $T(z) = 10 + 12z + 218.4\,(1 - e^{-z/10})$, $z$ em km, $T$ em
\unit{\degreeCelsius}.

\textbf{16.} $q(z) = 30 + 54.6\,e^{-z/10}$: \qty{50.1}{mW.m^{-2}} a \qty{10}{km},
\num{37.4} a \qty{20}{km} e \qty{31.7}{mW.m^{-2}} a \qty{35}{km}.

\textbf{17.} $30/84.6 = \qty{35}{\percent}$: quase dois terços do \omterm{def:b1:heat-budget:heat-flow}{fluxo de calor}
sobre E são produzidos em sua camada radiogênica.

\textbf{18.} $T(10) = 10 + 120 + 218.4\times0.632 = \qty{268}{\degreeCelsius}$.

\textbf{19.} $T(20) = 10 + 240 + 218.4\times0.865 = \qty{439}{\degreeCelsius}$.

\textbf{20.} Sob A, $A_0D^2/k = 1.35\times100/2.5 = \qty{54.0}{K}$:
$T(35) = 10 + 420 + 54.0\times0.970 = \qty{482}{\degreeCelsius}$.

\textbf{21.} Pelo $A_0D^2/k$: \qty{54}{K} sob A e \qty{218}{K} sob E. O fluxo profundo
de calor é o mesmo; a \omterm{def:b1:layered-earth:layers}{crosta} de E é quatro vezes mais rica em fontes de calor, e o calor
que elas produzem precisa ser conduzido para cima pela mesma rocha, o que exige um
gradiente mais íngreme em todo o trajeto.

\textbf{22.} $1300 = 10 + 218.4 + 12z$ dá $z = \qty{89.3}{km}$ sob E;
$1300 = 10 + 54.0 + 12z$ dá \qty{103.0}{km} sob A.

\textbf{23.} Ambas estão dentro da faixa da \omterm{def:b1:layered-earth:lithosphere}{litosfera} continental, de cem quilômetros ou menos sob regiões jovens até cerca de \qty{250}{km} sob
crátons, do lado fino dessa faixa: com um fluxo reduzido de \qty{30}{mW.m^{-2}} e esta
condutividade, o \omterm{def:b1:heat-budget:geotherm}{geotermo} condutivo atinge temperaturas do \omterm{def:b1:layered-earth:layers}{manto} a cerca de
\qty{100}{km}; sob o plúton mais radioativo, \qty{14}{km} antes.

\textbf{24.} $T(35) = 10 + 420 + 218.4\times(1 - e^{-3.5}) = 430 + 211.8$: a
\omterm{def:b1:layered-earth:discontinuity}{Moho} sob o plúton E está a
\textbf{\qty[text-series-to-math=true]{642}{\degreeCelsius}}.
\end{solution}
